\section{Bilan du Jalon 1 : Socle Technique Prêt pour les Contrats Métier}

\begin{tcolorbox}[colback=emerald!4!white,colframe=emerald!40!white,arc=2mm,boxrule=0.8pt,top=2mm,bottom=2mm,left=3mm,right=3mm]
    \small\textbf{\color{emerald!80!black}Synthèse d'Étape} : À l'issue du Jalon 1, l'ensemble des fondations d'ingénierie logicielle est posé. La solution est prête à accueillir les contrats de domaine du Jalon 2.
\end{tcolorbox}

\vspace{0.1cm}

Dans tout apprentissage rigoureux d'un système Fullstack, il est capital pour l'apprenant de faire un arrêt sur image pour mesurer avec exactitude **ce qui est d'ores et déjà opérationnel** et **ce qui relève des jalons ultérieurs**.

\vspace{0.2cm}

\subsection*{État Actuel du Système vs Objectifs Cibles}

\begin{table}[h!]
\centering
\small
\begin{tabularx}{\textwidth}{l p{5.5cm} X}
\toprule
\textbf{Composant} & \textbf{État Réalisé (Jalon 1)} & \textbf{Cible Planifiée (Jalons 2 à 15)} \\
\midrule
\textbf{\texttt{FullStackCSharp.Shared}} & Bibliothèque \texttt{net10.0} compilable & Hébergera tous les DTOs \texttt{record} (Jalon 2) \\
\textbf{\texttt{FullStackCSharp.Back}} & Kestrel autonome + Sonde \texttt{/health} & Pipeline HTTP, SQLite, EF Core 10, JWT (Jalons 3--7) \\
\textbf{\texttt{FullStackCSharp.Front}} & Application .NET MAUI compilable & MVVM, Shell, Refit, Panier local (Jalons 8--15) \\
\textbf{Configuration} & \texttt{Directory.Build.props} à la racine & Règle universelle de nullabilité active \\
\textbf{Liaisons} & Références \texttt{ProjectReference} câblées & Typage fort de bout en bout sur le réseau \\
\bottomrule
\end{tabularx}
\end{table}

\vspace{0.2cm}

\begin{schemabox}{Vue Consolidée de la Solution au Terme du Jalon 1}
\begin{center}
\begin{tikzpicture}[
    boxsln/.style={
        rectangle, 
        draw=accent!90!black, 
        fill=accent!6!white, 
        rounded corners=2.5mm, 
        minimum width=15.2cm, 
        minimum height=1.1cm, 
        inner ysep=2mm,
        align=center, 
        line width=1pt
    },
    boxproj/.style={
        rectangle, 
        draw=#1!80!black, 
        fill=#1!6!white, 
        rounded corners=2mm, 
        minimum width=4.7cm, 
        minimum height=2.0cm, 
        inner ysep=2.5mm,
        align=center, 
        line width=0.85pt
    },
    checkarrow/.style={
        -{Stealth[length=2.5mm]}, 
        thick, 
        color=emerald!80!black
    }
]

% Master Solution
\node[boxsln] (sln) at (0, 3.4) {
    \textbf{\large FullStackCSharp.sln} \textit{(Orchestrateur MSBuild unifié)}\\[1pt]
    \footnotesize Configuration globale : \texttt{Directory.Build.props} $\bullet$ Zéro avertissement $\bullet$ .NET 10 LTS
};

% 3 projects
\node[boxproj=cyanproj] (shared) at (-5.2, 0.8) {
    \textbf{\normalsize 1. Shared}\\[2pt]
    {\color{cyanproj}\rule{3.8cm}{0.4pt}}\\[2pt]
    \scriptsize Contrats d'échange\\
    \tiny Type : \texttt{classlib}\\
    \tiny \color{emerald!80!black}\textbf{\checkmark~Prêt pour les DTOs}
};

\node[boxproj=secondary] (back) at (0, 0.8) {
    \textbf{\normalsize 2. Back}\\[2pt]
    {\color{secondary}\rule{3.8cm}{0.4pt}}\\[2pt]
    \scriptsize Serveur HTTP Kestrel\\
    \tiny Route : \texttt{GET /health}\\
    \tiny \color{emerald!80!black}\textbf{\checkmark~Port 5227 validé}
};

\node[boxproj=indigo] (front) at (5.2, 0.8) {
    \textbf{\normalsize 3. Front}\\[2pt]
    {\color{indigo}\rule{3.8cm}{0.4pt}}\\[2pt]
    \scriptsize Client .NET MAUI\\
    \tiny Accueil : \textit{Gelato Artisan}\\
    \tiny \color{emerald!80!black}\textbf{\checkmark~Moteur XAML validé}
};

\draw[checkarrow] (-5.2, 2.8) -- (shared.north) node[midway, left=2pt, font=\tiny\bfseries\color{emerald!80!black}] {Compilé};
\draw[checkarrow] (0, 2.8) -- (back.north) node[midway, right=2pt, font=\tiny\bfseries\color{emerald!80!black}] {Compilé};
\draw[checkarrow] (5.2, 2.8) -- (front.north) node[midway, right=2pt, font=\tiny\bfseries\color{emerald!80!black}] {Compilé};

\end{tikzpicture}
\end{center}
\end{schemabox}

\vspace{0.2cm}

\begin{bestpracticebox}{Check-list de Validation Finale du Jalon 1}
Avant d'aborder le Jalon 2, l'apprenant doit pouvoir valider les 3 points suivants :
\begin{enumerate}
    \item La commande \texttt{dotnet build FullStackCSharp.sln} s'exécute avec succès en retournant \textbf{0 avertissement et 0 erreur}.
    \item L'exécution de \texttt{dotnet run --project FullStackCSharp.Back} ouvre le serveur sur le port \texttt{5227} et l'appel de \texttt{/health} retourne le statut JSON \texttt{"Healthy"}.
    \item Les fichiers \texttt{.csproj} du Back et du Front intègrent chacun leur référence vers \texttt{Shared}, préfigurant l'autorité centrale des types partagés.
\end{enumerate}
\end{bestpracticebox}
